\begin{lemma}[Transitivity of the lifted relation]
If $r$ is a preorder on $Q$, then for all
$X,Y,Z\in\mathsf{PowerQ}(Q)$,
\[
\mathsf{PowerRel}(r)(X,Y)\land\mathsf{PowerRel}(r)(Y,Z)
\Longrightarrow\mathsf{PowerRel}(r)(X,Z).
\].
\end{lemma}
\begin{proof}
Use well-founded induction on the lexicographic pair of child relations used
in \noderef{PowerRel}; the required well-foundedness is obtained from
\noderef{PowerChildWellFounded}.  Split $X,Y,Z$ by their atom or node tags
and use the four reductions \noderef{PowerRelAtomAtom},
\noderef{PowerRelAtomNode}, \noderef{PowerRelNodeAtom}, and
\noderef{PowerRelNodeNode}.

For three atoms, transitivity is exactly transitivity of $r$.  If the middle
term is a node, the first comparison supplies a child of that node whenever
the left term is an atom, while the second comparison supplies a child of
the right term for that middle child; compose the two recursive comparisons.
If the left term is a node, perform this argument separately for each left
child.  The node--atom case instead universally supplies comparisons from
each left child into the middle atom, which then composes with the second
comparison.  In the node--node case, for each child of the first node first
choose the corresponding child of the middle node and then the corresponding
child of the final node; the induction hypothesis composes the two child
comparisons.  These are exactly the existential and universal quantifier
orders in the cited reductions, and all recursive calls descend in the
lexicographic child relation.  Thus the resulting witnesses establish
$\mathsf{PowerRel}(r)(X,Z)$ in every tag case.
\end{proof}
